\subsection{正合列}

定义 5. 设 F, G, H 为三个 A-模；设 f 是 F 到 G 中的同态，g 是 G 到 H 中的同态。有序对 \((f, g)\) 称为正合列，如果

\[
\bar {g} ^ {- 1} (0) = f (\mathbf {F})
\]

换言之，即 g 的核等于 f 的像。

该图表

\[
\mathrm{F} \xrightarrow {f} \mathrm{G} \xrightarrow {g} \mathrm{H}\tag{12}
\]

也称为正合列。

我们类似地考虑由四个 A-模和三个同态构成的图表：

\[
\mathrm{E} \xrightarrow {f} \mathrm{F} \xrightarrow {g} \mathrm{G} \xrightarrow {h} \mathrm{H}.\tag{13}
\]

若图表 \(E \xrightarrow{f} F \xrightarrow{g} G\) 是正合的，则称此图表在 F 处正合；若 \(F \xrightarrow{g} G \xrightarrow{h} H\) 是正合的，则称它在 G 处正合。若图表 (13) 在 F 处和 G 处都正合，则简称它是正合的，或称为正合列。具有任意多个项的正合列可类似地定义。

注记。(1) 若有序对 \((f, g)\) 是一个正合列，则 \(g \circ f = 0\) ；但这一性质当然并不刻画正合列，因为它只表明 f 的像包含在 g 的核中。

在下面的陈述中，E, F, G 表示 A-模，0 表示约化为其单位元的 A-模；箭头表示 A-模同态。由于从模 0 到模 E（相应地，从 E 到 0）只有一个同态，故在出现这些同态的正合列中无需给它们命名。

命题 3. (a) 要使

\[
0 \longrightarrow \mathbf {E} \stackrel {f} {\longrightarrow} \mathbf {F}
\]

成为正合列，必须且只需 f 是单射。

\begin{enumerate}
\def\labelenumi{(\alph{enumi})}
\setcounter{enumi}{1}
\tightlist
\item
  要使
\end{enumerate}

\[
\mathbf {E} \xrightarrow {f} \mathbf {F} \longrightarrow 0
\]

成为正合列，必须且只需 \(f\) 是满射。

\begin{enumerate}
\def\labelenumi{(\alph{enumi})}
\setcounter{enumi}{2}
\tightlist
\item
  要使
\end{enumerate}

\[
0 \longrightarrow \mathrm{E} \stackrel {f} {\longrightarrow} \mathrm{F} \longrightarrow 0
\]

成为正合列，必须且只需 \(f\) 是双射（换言之（第 2 号，命题 2），即 \(f\) 是 \(\mathbf{E}\) 到 \(\mathbf{F}\) 上的同构）。

\begin{enumerate}
\def\labelenumi{(\alph{enumi})}
\setcounter{enumi}{3}
\tightlist
\item
  若 \(\mathbf{F}\) 是 \(\mathbf{E}\) 的一个子模， \(i: \mathbf{F} \to \mathbf{E}\) 是典范单射， \(p: \mathbf{E} \to \mathbf{E} / \mathbf{F}\) 是典范同态，则图表
\end{enumerate}

\[
0 \longrightarrow \mathrm{F} \stackrel {i} {\longrightarrow} \mathrm{E} \stackrel {p} {\longrightarrow} \mathrm{E} / \mathrm{F} \longrightarrow 0\tag{14}
\]

是一个正合列。

\begin{enumerate}
\def\labelenumi{(\alph{enumi})}
\setcounter{enumi}{4}
\tightlist
\item
  若 \(f: \mathbf{E} \to \mathbf{F}\) 是一个同态，则图表
\end{enumerate}

\[
0 \longrightarrow f ^ {- 1} (0) \stackrel {i} {\longrightarrow} \mathrm{E} \stackrel {f} {\longrightarrow} \mathrm{F} \stackrel {p} {\longrightarrow} \mathrm{F} / f (\mathrm{E}) \rightarrow 0
\]

（其中 i 是典范单射，p 是典范满射）是一个正合列。

该命题直接由定义及第 2 号的命题 2 得出。

注记。(2) 说存在一个正合序列

\[
0 \longrightarrow \mathrm{E} \stackrel {f} {\longrightarrow} \mathrm{F} \stackrel {g} {\longrightarrow} \mathrm{G} \longrightarrow 0
\]

意味着 \(f\) 是单射， \(g\) 是满射，并且与 \(g\) 相伴的典范双射是 \(\mathrm{F} / f(\mathbf{E})\) 到 \(\mathbf{G}\) 上的一个同构。也称三元组 \((\mathbf{F}, f, g)\) 是模 \(\mathbf{G}\) 借助模 \(\mathbf{E}\) 的扩张（I，§6，第 7 号）。

\begin{enumerate}
\def\labelenumi{(\arabic{enumi})}
\setcounter{enumi}{2}
\tightlist
\item
  若存在一个 4 项的正合列
\end{enumerate}

\[
\mathrm{E} \xrightarrow {f} \mathrm{F} \xrightarrow {g} \mathrm{G} \xrightarrow {h} \mathrm{H}
\]

则 \(f\) 的余核是 \(\mathrm{F} / f(\mathrm{E}) = \mathrm{F} / g^{-1}(0)\) ，而 \(h\) 的核是 \(g(\mathrm{F})\) ；因此与 \(g\) 相伴的典范双射是一个同构

\[
\operatorname{Coker} f \cong \operatorname{Ker} h.
\]

\begin{enumerate}
\def\labelenumi{(\arabic{enumi})}
\setcounter{enumi}{3}
\tightlist
\item
  考虑一个 A-模同态的有序对
\end{enumerate}

\[
\mathrm{E} \xrightarrow {f} \mathrm{F} \xrightarrow {g} \mathrm{G}.\tag{15}
\]

要使图表 (15) 成为正合列，必须且只需存在两个 A-模 S, T 以及同态 \(a:E\to S\) , \(b:S\to F\) , \(c:F\to T\) , \(d:T\to G\) ，使得下列三个序列

\[
\left\{ \begin{array}{c} \mathrm{E} \xrightarrow {a} \mathrm{S} \longrightarrow 0 \\ 0 \longrightarrow \mathrm{S} \xrightarrow {b} \mathrm{F} \xrightarrow {c} \mathrm{T} \longrightarrow 0 \\ 0 \longrightarrow \mathrm{T} \xrightarrow {d} \mathrm{G} \end{array} \right.\tag{16}
\]

都是正合的，并且 \(f = b \circ a\) ， \(g = d \circ c\) 。

若 (15) 是正合列，则取 \(\mathbf{S} = f(\mathbf{E}) = \overline{g}^{-1}(0)\) 与 \(\mathbf{T} = g(\mathbf{F})\) ，其中 \(b\) 和 \(d\) 是典范单射，而 \(a\) （相应地 \(c\) ）是与 \(f\) （相应地 \(g\) ）有相同图像的同态。反之，若 \(\mathbf{S}, \mathbf{T}, a, b, c, d\) 满足上述条件，则 \(f(\mathbf{E}) = b(a(\mathbf{E})) = b(\mathbf{S})\) ， \(g^{-1}(0) = c^{-1}(d^{-1}(0)) = c^{-1}(0)\) ，因此 (16) 的正合性表明 \(f(\mathbf{E}) = \overline{g}^{-1}(0)\) 。

在正合列中，当用显式字母表示同态对论证并非必要时，往往可以省略这些字母。

注记。(5) 正合列的定义可立即推广到未必交换的群；此时当然采用乘法记号，并在公式中将 0 换成 1（若不致引起混淆）。命题 3 的 (a)、(b)、(c) 仍然成立，当 F 是 E 的正规子群时 (d) 也成立。注记 2 和命题 3 (e) 在 \(f(\mathrm{E})\) 是 F 的正规子群的条件下成立；注记 3 和注记 4 无需修改即有效。

